%% ===================================================================
\section{Introduction}
\label{sec:intro}
%% ===================================================================

Let $G$ be a finite group and let $\mathcal C\subseteq G$ be a
conjugation-closed set generating $G$, for instance a single conjugacy
class. Label the $2^m$ leaves of a balanced binary tree of depth $m$ by
elements of $\mathcal C$ and let every internal vertex take the commutator
$[x,y]=xyx^{-1}y^{-1}$ of its two children. Expanding the root value as a
word in the leaf labels gives a product of $4^m$ occurrences, each a leaf
label or its inverse. Now read every occurrence independently: with
probability $1/2$ it contributes its element, otherwise the identity. The
output is a random product of $4^m$ independent lazy factors, that is, a
lazy random walk on $G$ whose steps are prescribed by the tree rather than
chosen freely. We ask when every such walk is uniformly close to the
uniform distribution, and what determines the depth at which this begins.

Two features distinguish the construction from a word map. First,
repeated occurrences of a label, including the inverse copies inside each
commutator, receive independent coin tosses, so the output is not the
value of a word at random arguments; in particular it is not confined to
the derived subgroup, and one-dimensional characters are contracted by
the individual lazy factors. Second, only \emph{valid} configurations are
considered, those whose deterministic commutator value at the root is
non-trivial. This constraint is what brings the derived series into the
problem: a configuration whose leaves all lie in a solvable subgroup $H$
has trivial root value as soon as the depth reaches the derived length of
$H$.

In the language of random walks, \emph{strict contraction} of a
configuration means that its Fourier operator has norm less than one in
every non-trivial irreducible representation, that is, a spectral gap on
the orthogonal complement of the constants; a \emph{uniform
strict-contraction level} is a depth at which every valid configuration
contracts strictly. When the labels at the $2^m$ leaf positions may be
chosen independently from $\mathcal C$ we speak of the \emph{configurable}
model, and the independent coin tosses at the $4^m$ expanded occurrences
constitute the \emph{independent-read} model; precise definitions are in
\Cref{sec:prelim}. A subgroup is \emph{$\mathcal C$-generated} if it is
generated by its intersection with $\mathcal C$.

\subsection*{The solvability dichotomy}
The answer is controlled by solvability, and the depth is measured by
derived length.

\begin{lettertheorem}[Solvability dichotomy and exact level]
\label{ithm:A}
Let $G$ be a finite non-solvable group and $\mathcal{C}$ a configurable conjugation-closed
leaf set with $\langle\mathcal{C}\rangle=G$. The construction has a uniform
strict-contraction level if and only if every proper $\mathcal{C}$-generated subgroup of
$G$ is solvable; if some proper $\mathcal{C}$-generated subgroup is non-solvable, the
obstruction is permanent: every depth admits a valid configuration that fails to contract in some non-trivial irreducible. When the level exists, it
is exactly
$D=\max\{\operatorname{dl}(H):H<G,\ H\text{ is }\mathcal{C}\text{-generated}\}$, the
largest derived length among proper $\mathcal{C}$-generated subgroups
\textup{(\Cref{thm:dichotomy})}.
\end{lettertheorem}

Equivalently, \Cref{ithm:A} determines the exact depth at which a
non-identity iterated commutator certifies that its entries generate the
whole group. Transpositions illustrate the two sides. In $S_5$ every
proper subgroup generated by transpositions is solvable, the largest
derived length is $3$, and the level is $3$. In $S_n$ for $n\ge6$ the
point stabiliser $S_{n-1}$ is a proper non-solvable subgroup generated by
transpositions: at every depth it supplies valid configurations supported
in a proper subgroup, and uniform contraction fails (\Cref{cor:Sn}).

The criterion characterises a classical family. A finite non-abelian
simple group is \emph{minimal simple} in Thompson's
sense~\cite{Thompson1968} if all of its proper subgroups are solvable, and
\Cref{ithm:A} applied to every class at once gives: a finite non-abelian
simple group is minimal simple if and only if every non-trivial conjugacy
class has a uniform strict-contraction level (\Cref{cor:minsimple}). In
that case every class has level at most $5$, by the bound on derived
lengths of proper subgroups of minimal simple
groups~\cite[Remark 2.3]{ErsoyTortoraTota2014}, and the involution class
of $\PSL_3(3)$ attains $5$ (\Cref{prop:psl33-level}). For involution
leaves in $A_5$ the level is $2$: although $A_4$ is a proper subgroup of
derived length $2$, its involutions generate only $V_4$, so only
$\mathcal C$-generated subgroups enter the criterion (\Cref{cor:A5}).

The proof of \Cref{ithm:A} rests on two ingredients. The analytic one
concerns a single configuration. For a tree $\tau$, let $\mu_\tau$ be its
unconditional output law and write $M_\tau(\rho)=\widehat{\mu_\tau}(\rho)$
for its Fourier operator; $V_\rho^H$ denotes the subspace fixed by $H$,
and operator norms are written $\norm{\cdot}$.

\begin{lettertheorem}[Fourier fixed-vector criterion]
\label{ithm:B}
For any unitary representation $\rho$ and any configuration $\tau$ with leaf set $S$ and
$H=\langle S\rangle$, one has $\opn{M_\tau(\rho)}=1\iff V_\rho^H\neq0$. Hence the
configuration contracts strictly in every non-trivial irreducible if and only if
$\langle S\rangle=G$: strict contraction $\iff$ generation $\iff$ escape from every proper
subgroup \textup{(\Cref{lem:geom,thm:criterion})}.
\end{lettertheorem}

The algebraic ingredient is the normal-subset word-value lemma of
Fern\'andez-Alcober and Pintonello~\cite[Lemma 2.1(ii)]{FernandezPintonello2023}:
for a conjugation-closed subset of a subgroup $H$, the depth-$k$
commutator values of its elements generate $H^{(k)}$. It identifies the
subgroup generated by realisable depth-$k$ targets with $H^{(k)}$, so the
validity constraint forces escape from every proper solvable
$\mathcal C$-generated subgroup at depth $D$, and from none of the
non-solvable ones; \Cref{sec:notfolklore} summarises the correspondence.
Testing all non-trivial irreducibles matters: a proper subgroup can act
transitively on a chosen permutation set, so contraction on its
augmentation representation alone need not detect the obstruction
(\Cref{prop:std}; $S_5$ on six points gives explicit examples).

\subsection*{A sharp angle between two octahedral subgroups}
Once a common fixed vector is absent, the angle between the fixed spaces
of two subgroups measures how strongly the product of their averaging
projections contracts; for conjugate subgroups this angle is the norm of
a finite Hecke compression (\Cref{cor:double-coset-angle}). The main
concrete theorem of the paper computes this angle, sharply and uniformly,
for a natural family in rank one.

For an odd prime $q$ the group $\PGL_2(q)$ contains subgroups isomorphic
to $S_4$; for $q\equiv\pm3\pmod 8$ they are \emph{exterior}, that is, not
contained in $\PSL_2(q)$, and they form a single conjugacy
class~\cite{Huppert1967,King2005,Giudici2007}. Reducing Hamilton's
quaternions modulo $q$ produces an explicit such subgroup $A_q$ together
with a conjugate $B_q=hA_qh^{-1}$, $h=i+j+k$, meeting $A_q$ in an $S_3$.
We prove that this pair is the only one of its kind, and that its angle is
governed by the Ramanujan graphs of
Lubotzky--Phillips--Sarnak~\cite{LPS1988}.

\begin{lettertheorem}[The natural octahedral pair]
\label{ithm:F}
For primes $q\equiv\pm3\pmod8$, $q\ne3$, the Hamilton construction
defines exterior subgroups $A_q,B_q\cong S_4$ of $G_q=\PGL_2(q)$ with
$A_q\cap B_q\cong S_3$ and $\langle A_q,B_q\rangle=G_q$. Every pair
$A_1,A_2$ of exterior octahedral subgroups of $G_q$ with $|A_1\cap A_2|=6$
is simultaneously conjugate to $(A_q,B_q)$, and its fixed-space
projections satisfy
\[
 \max_{\rho\in\widehat{G_q},\,\rho\ne\mathbf1}
              \opn{P_{A_1}^\rho P_{A_2}^\rho}=c_q\le\frac{\sqrt3}{2},
 \qquad c_q\longrightarrow\frac{\sqrt3}{2},
\]
where $c_q$ is the common value attained by the natural pair
\textup{(\Cref{prop:octahedral-cap6,thm:octahedral-sharp})}.
\end{lettertheorem}

The mechanism is a compression identity: the four norm-$3$ quaternions
$\pm i\pm j\pm k$ represent the double coset $A_qhA_q$, so the Hecke
operator of the $4$-regular Lubotzky--Phillips--Sarnak graph $X^{3,q}$,
restricted to the $A_q$-fixed vectors of each irreducible, \emph{is} the
angle operator $P_{A_q}\rho(h)P_{A_q}$. The Ramanujan bound $2\sqrt3$,
normalised by the degree $4$, is $\sqrt3/2=\cos(\pi/6)$; this is the
origin of the constant. Alon--Boppana on the quotient graphs gives
optimality. Geometrically, $\theta_q=\arccos c_q$ is the smallest
fixed-space angle over all non-trivial irreducibles, $\theta_q\ge\pi/6$
and $\theta_q\to\pi/6$. Combined with an internal $S_4$ estimate, the
angle bound yields an explicit uniform contraction constant for the
corresponding tree blocks (\Cref{cor:octahedral-block}).

\Cref{ithm:F} settles the intersection-six stratum of the following
question, raised in the preliminary rank-one
preprint~\cite[Conjecture 6.9]{ZermenoRankOneZenodo}.

\begin{conjecture}[Uniform angle for exterior octahedral pairs]
\label{conj:allpairs}
For primes $q\equiv\pm3\pmod8$, $q\ne3$,
\[
 \sup_{\substack{A,B\cong S_4\ \mathrm{exterior}\\
                 \langle A,B\rangle=\PGL_2(q)}}\ \
 \max_{\rho\ne\mathbf1}\opn{P_A^\rho P_B^\rho}
 \le\frac{\sqrt3}{2}+o(1)\qquad(q\to\infty).
\]
\end{conjecture}

By \Cref{prop:octahedral-cap6} two distinct exterior octahedral subgroups
meet in a subgroup of order at most $6$, so the remaining cases are the
intersection orders $1,2,3,4$. \Cref{tab:allpairs-evidence} records numerical
all-pairs maxima for five primes up to $37$; all lie below $\sqrt3/2$.

\subsection*{Classifications}
The subgroup criterion yields exact levels in concrete families. In
projective rank one it gives a full involution-class classification. An
\emph{exterior} involution of $\PGL_2(q)$ is one whose determinant is a
nonsquare; the other involutions lie in $\PSL_2(q)$.

\begin{lettertheorem}[Rank-one involutions]
\label{ithm:D}
For odd $q\ge5$, all involutions of $\PSL_2(q)$ are obstructed, whereas the
exterior class is obstructed exactly by proper subfields of size at least $5$
and odd extension degree. In the unobstructed exterior cases the exact level
is $3$ for $q\equiv\pm3\pmod8$ and $2$ for $q\equiv\pm1\pmod8$.
The involution classes of $\PSL_2(q)$, including its small exceptions, are
classified in \Cref{thm:rankone-classification}.
\end{lettertheorem}

The dependence on the leaf set is visible in $\PSL_3(3)$: its involution
class has level $5$ while its transvection class has level $4$.
\Cref{prop:transvection-levels} classifies transvection leaves in
$\SL_n(q)$ for $n\ge3$: the level is $3$ for $\SL_3(2)$, $4$ for
$\SL_3(3)$, and in every other case a proper non-solvable subgroup
generated by transvections obstructs permanently. For a generating
conjugacy class of $\GL_2(\Z/p^r)$, $p\ge5$, \Cref{prop:congruence-class}
transfers the dichotomy from the reduction modulo $p$; at $r=2$ the level
increases by at most one.

For \emph{rigid} templates, one fixed leaf per position as in an actual
branching program, the configurable level $D$ becomes an upper bound, and
persistence of validity is a finite dynamical question on conjugacy
classes, which we solve.

\begin{lettertheorem}[Rigid templates]
\label{ithm:C}
Let $\mathcal{C}$ be conjugation-closed.
\textup{(i)} For a non-empty set $S$ of non-trivial conjugacy classes, valid binary
compositions persist at every level from $S$ if and only if some class in $S$ has
non-solvable normal closure---equivalently, escapes the solvable radical
$O_\infty(G)$---if and only if some class in $S$ reaches a cycle of the target transition
graph $\mathcal{T}(G)$; and mixed-class children never extend the reach of persistence
beyond conjugated doubling \textup{(\Cref{thm:diagred,cor:binaryred,cor:binarylevel})}.
\textup{(ii)} Over the transposition class of $S_5$, the non-vacuously realisable exact
levels of rigid families are precisely $\{2,3\}$ \textup{(\Cref{thm:s5levels})}.
\end{lettertheorem}

\subsection*{Quantitative consequences}
For a fixed group, strict contraction gives a positive gap by finiteness;
over a growing family the gap may tend to zero. For involution leaves and
uniformly bounded depth, a uniform contraction gap is equivalent to a
uniform spectral gap for the associated lazy Cayley walks, each of which
chooses an expanded leaf occurrence uniformly and applies it with
probability $1/2$; \Cref{prop:cayley-comparison,cor:bounded-depth-gap}
give explicit two-sided comparisons, so the block problem becomes an
expansion problem for its particular generating sets.

The next result controls the mixing depth even when the base gap decreases
with group size. Write $U_G$ for the uniform law on $G$.

\begin{lettertheorem}[Uniform rates and optimal order of depth]
\label{ithm:E}
For the unobstructed classes of \Cref{ithm:D} and
\Cref{thm:rankone-classification}, and for every non-trivial conjugacy class
of every minimal simple group, there are absolute $c,B>0$ such that every
valid tree $\tau$ of level $m\ge D$ satisfies
\[
 \dtv(\mu_\tau,U_G)\le\tfrac12\sqrt{|G|-1}
                  \exp\bigl(-c4^m/(\log|G|)^B\bigr).
\]
For fixed error in $(0,1/2)$ the least depth guaranteeing this error for
all valid trees has order $\Theta(\log\log|G|)$ as $|G|\to\infty$
\textup{(\Cref{thm:family-rate,prop:depth-lower,cor:family-examples})}.
\end{lettertheorem}

The proof converts the bounded-rank diameter theorem of
Pyber--Szab\'o~\cite{PyberSzabo2016} into a base-block contraction estimate
and amplifies it through the tree. The lower bound is elementary: at depth
$m$ the output is supported on at most $2^{4^m}$ elements, which forces
depth $\Omega(\log\log|G|)$ for fixed-error mixing. Over prime fields
outside a set of relative density zero among the primes, the strong
uniform expansion theorem of Breuillard--Gamburd~\cite{BreuillardGamburd2010}
gives a constant gap for every generating involutive block of bounded
depth, which for the unobstructed classes improves the exponential factor
to $e^{-c4^m}$ (\Cref{cor:almost-all-primes}). Over arbitrary finite
fields, the dimension-zero expansion theorem of
Becker--Breuillard~\cite{BeckerBreuillard2026} bounds the number of
exceptional simultaneous-conjugacy classes of generating involutive
configurations at each fixed depth by $O(q^\delta)$ for every $\delta>0$
(\Cref{cor:few-exceptional-blocks}).

\subsection*{Relation to the literature}
\medskip\noindent\textbf{Probabilistic generation.}
That random elements of a finite simple group generate it with
probability tending to one is the theorem of
Dixon~\cite{Dixon1969}, Kantor--Lubotzky~\cite{KantorLubotzky1990} and
Liebeck--Shalev~\cite{LiebeckShalev1995}; the survey of
Burness~\cite{Burness2019} describes the subsequent theory of generation
by elements of prescribed type, including the involution results
of~\cite{LiebeckShalev1996}, and Harper--Quick~\cite{HarperQuick2026}
have recently extended the asymptotic theorem beyond simple groups.
\Cref{ithm:B} turns a generation question into a spectral one:
contraction holds exactly when the leaves generate, and \Cref{ithm:A}
identifies the proper subgroups that can obstruct it as the non-solvable
$\mathcal C$-generated ones. Random generation asks how often a tuple
escapes every proper subgroup; here the validity constraint forces the
escape, and derived length records when.

\medskip\noindent\textbf{Word maps and commutator distributions.}
Word maps connect random inputs to the distribution of word
values~\cite{LarsenShalev,Shalev2013}. By the Ore
conjecture~\cite{LOST2010} every element of a finite non-abelian simple
group is a commutator, and Garion--Shalev~\cite{GarionShalev2009} prove
that iterated commutators with any bracketing, evaluated at uniform
arguments, are asymptotically equidistributed on finite simple groups of
growing order. Schul--Shalev~\cite{SchulShalev2011} show that the random
walk whose steps are uniform on the values of a fixed non-trivial word
mixes in a bounded number of steps, and
Larsen--Shalev--Tiep~\cite{LarsenShalevTiep2019} obtain almost uniform
distributions for products of word values in disjoint variables. These
results concern a fixed word at uniform arguments on simple groups of
large order. The present construction differs in three respects: the
arguments are restricted to a conjugation-closed set, every expanded
occurrence is read independently (the operator is built from
$XY\tilde X^{-1}\tilde Y^{-1}$ with independent copies,
\Cref{lem:tree,rem:notxy}), and the group is fixed while the depth varies.
The questions are accordingly different: not whether a word
equidistributes asymptotically, but which subgroup configurations
obstruct contraction and at what depth the obstruction disappears.

\medskip\noindent\textbf{Random walks driven by conjugacy classes.}
When the leaf set is a single class $C$, averaging over $C$ has
eigenvalues $\chi(C)/\chi(1)$, and the mixing of class-driven walks is a
classical subject: random transpositions~\cite{DS81}, random
transvections~\cite{Hildebrand1992}, prescribed cycle
types~\cite{LulovPak,Roichman1996}, the character-theoretic estimates of
Liebeck--Shalev for groups of Lie type~\cite{LiebeckShalev2005}, and the
mixing-time-two phenomenon for the class of a random
element~\cite{Shalev2008}; the representation-theoretic framework is
Diaconis's~\cite{Diaconis1988}. In the commutator tree the leaf-subgroup
criterion detects trapping through fixed spaces in $\Ind_H^G\mathbf1$
rather than through character bounds, and the fixed-group estimate in
\Cref{sec:prelim} uses the Diaconis--Shahshahani bound~\cite{DS81}.

\medskip\noindent\textbf{Solvability detected by small subgroups.}
The appearance of solvability through proper subgroups generated by
prescribed elements has several precedents. Thompson's
classification~\cite{Thompson1968} characterises minimal simple groups
by the solvability of their proper subgroups; the solvable radical is the
set of elements $y$ with $\langle x,y\rangle$ solvable for all
$x$~\cite{GKPS2006}, and admits descriptions through boundedly many
conjugates and commutators~\cite{GGKP2008}; finite solvable groups are
characterised by recursively defined two-variable
laws~\cite{Bandman2006,BrayWilsonWilson2005}; and a finitely generated
linear group satisfies a probabilistic identity exactly when it is
virtually solvable~\cite{LarsenShalev2016}. \Cref{ithm:A} belongs to this
circle of ideas with an analytic conclusion: solvability of the proper
$\mathcal C$-generated subgroups is equivalent to the eventual onset of a
uniform spectral gap, and their derived lengths give the exact onset.

\medskip\noindent\textbf{Expansion, quasirandomness and angles.}
Gowers~\cite{Gowers2008} and Babai--Nikolov--Pyber~\cite{BabaiNikolovPyber2008}
show that a group with no small non-trivial representations mixes
products of large sets; this quasirandomness is an input to the
Bourgain--Gamburd method, which with product
growth~\cite{Helfgott2008,BGT2011,PyberSzabo2016} yields uniform spectral
gaps for Cayley graphs of $\SL_2(\F_p)$~\cite{BourgainGamburd},
of its prime-power quotients~\cite{BourgainGamburdJEMS2008,BourgainVarju2012},
of perfect groups over square-free moduli~\cite{SalehiVarju}, and for
random generators in bounded rank~\cite{BGGT2015}; finite simple groups
admit uniformly expanding generating sets~\cite{KLN2006}, while bounded
derived length obstructs expansion in solvable
families~\cite{LubotzkyWeiss1993}. In the tree model a common fixed
vector for a proper leaf subgroup is the exact obstruction, and
\Cref{cor:bounded-depth-gap} relates a uniform gap at bounded depth to
expansion of the particular leaf walks, using Gao's projection union bound
as presented in~\cite{AnshuAradVidick2016}. The passage from fixed spaces
to quantitative contraction is the theory of angles between subspaces and
alternating projections~\cite{KayalarWeinert1988,BadeaGrivauxMuller2010},
which Kassabov~\cite{Kassabov2011} uses to obtain Kazhdan constants from
arrangements of subgroup-fixed spaces; \Cref{ithm:F} computes one such
angle exactly from arithmetic.

\medskip\noindent\textbf{Freely chosen random products.}
Amir, Blachar, Ghosh and Vishne~\cite{AmirMixability} study
\emph{mixability}: the existence of a finite random product that is
exactly uniform, with its steps and Bernoulli probabilities freely chosen;
every $S_n$ is mixable, and Angel--Holroyd~\cite{AngelHolroyd2018}
construct exact shuffles from lazy adjacent transpositions. Their
representation criterion $\prod_i((1-p_i)I+p_i\rho(g_i))=0$ has a
singular-step lemma that is a linear cousin of \Cref{lem:geom}. Our
question instead fixes fair probabilities and a commutator pattern, and
asks for contraction uniformly over all valid configurations: among $S_n$
with $n\ge5$ and transposition leaves, only $S_5$ has such a level.

\subsection*{Origin and applications}
The construction is Barrington's~\cite{Bar89}: width-five branching
programs over $S_5$ compute $\mathrm{NC}^1$ by combining their instructions
through a balanced tree of commutators, and over matrix groups one may
use pairs conjugate to their commutator~\cite{Boneh2013}. Group programs
of this kind are randomised in secure computation~\cite{Kilian1988}, and
the distribution of products of random group elements underlies
leakage-resilient circuits built from $A_5$~\cite{MilesViola2013} and the
interleaved-product bounds in $\SL_2(q)$ of
Gowers--Viola~\cite{GowersViola2019}. The question answered here is what
happens when each expanded instruction is independently replaced by the
identity with probability $1/2$. The companion
paper~\cite{ZermenoS5} proves target-independence and a quantitative
mixing bound for the $S_5$ construction on inconsistent reads; the
present paper establishes the general structural law, its classifications
and its quantitative consequences. For the conditioned law,
\Cref{thm:tindep,lem:conditioning} give target-independence and an
additive conditioning error.

\subsection*{Reading guide}
After the model and the conditioned branching-program law in \Cref{sec:prelim},
readers interested in the structural results can proceed to
\Cref{sec:criterion,sec:notfolklore} for \Cref{ithm:A,ithm:B}, and to
\Cref{subsec:rigid} for \Cref{ithm:C}.
\Cref{sec:layers,sec:synthesis,sec:rankone} give obstruction mechanisms,
examples and the classification of \Cref{ithm:D}.
The quantitative consequences and the octahedral pair are developed in
\Cref{sec:quantitative-families,sec:octahedral-pair}.
\Cref{sec:computation} distinguishes exact certificates from numerical
checks, and \Cref{sec:open} records the remaining problems.
